\documentclass[11pt]{amsart}
\usepackage[margin=1in]{geometry}
\usepackage[T1]{fontenc}
\usepackage{lmodern}
\usepackage{amsmath,amssymb,mathtools}
\usepackage{booktabs}
\usepackage{microtype}
\usepackage[hidelinks]{hyperref}
\hypersetup{pdftitle={A Parseval-Prefix Improvement for Erdos Minimum-Overlap Problem},pdfauthor={Dmitry Khanukov},pdfsubject={Preliminary computer-assisted lower bound for Erdos problem 36},pdfkeywords={Erdos, minimum overlap, Parseval, computer-assisted proof}}
\usepackage{url}

\newtheorem{theorem}{Theorem}
\newtheorem{lemma}[theorem]{Lemma}
\newtheorem{corollary}[theorem]{Corollary}
\newtheorem{remark}[theorem]{Remark}
\DeclareMathOperator{\sinc}{sinc}
\newcommand{\R}{\mathbb R}
\newcommand{\Z}{\mathbb Z}
\newcommand{\dd}{\,d}
\newcommand{\one}{\mathbf 1}

\title{A Parseval-Prefix Improvement for Erd\H{o}s' Minimum-Overlap Problem}
\input{author-config.tex}
\date{August 16, 2026}

\begin{document}

\begin{abstract}
We give a computer-assisted proof that the asymptotic constant in Erd\H{o}s'
minimum-overlap problem satisfies
\[
  c_E>0.3805603.
\]
The previous certified lower bound was
$0.380554702762594012\ldots$.  The new ingredient is a finite
Parseval-prefix inequality, projected directly into the overlap-profile dual
language and used to replace the two binding central mean bins of a public
172-bin certificate.  A directed-rounding MPFR program verifies the new
central positive-part integral; the remaining 170 bins are inherited from
pinned Arb reports with verified hashes.  The optimizer that found the
multipliers is not part of the proof.  This manuscript is preliminary and
unrefereed, and the proof has not been formalized or verified in Lean.
\end{abstract}

\maketitle

\begin{center}
\textbf{Status: Preliminary and unrefereed.}
\end{center}

\section{Introduction}

For a partition
\[
 A\sqcup B=\{1,\ldots,2N\},\qquad |A|=|B|=N,
\]
let
\[
 r_{A,B}(m)=|\{(a,b)\in A\times B:a-b=m\}|,
 \qquad
 M(N)=\min_{A,B}\max_m r_{A,B}(m).
\]
Define
\[
 c_E:=\liminf_{N\to\infty}\frac{M(N)}{N}.
\]
The minimum-overlap problem, originating with Erd\H{o}s
\cite{Erdos1955,Moser1959}, asks for the asymptotic behavior of $M(N)$.

Haugland developed strong upper constructions
\cite{Haugland1996,Haugland2016}.  White introduced the modern Fourier and
convex-optimization approach to lower bounds \cite{White2022}.  Price later
publicly posted a compact positive-part dual certificate, checked with Arb, proving
$c_E>0.38055470$ and recording the sharper numerical consequence
$0.380554702762594012\ldots$ \cite{Price2026}.

The purpose of this note is to improve that certified lower bound.  Parseval
constraints were already present in White's larger optimization framework.
Our contribution is not the discovery of Parseval's identity; it is the direct
profile-space prefix inequality below, its use in a compact dual certificate,
and a new directed-rounding verification.

The source package and reproduction instructions are available at
\begin{center}
\url{https://github.com/khanukov/erdos36}.
\end{center}

\subsection*{Status and trust boundary}
This is a preliminary, unrefereed computer-assisted result.  The two central
mean bins are checked by the new directed-rounding MPFR verifier included in
the repository.  The other 170 bins are not rerun locally with Arb by this
package: they are inherited from Price's publicly posted per-bin Arb reports at the
pinned upstream commit, downloaded and checked against recorded SHA-256
hashes, after which the package recomputes their largest reported upper
endpoint.  Thus the outer-bin part trusts the pinned Price certificate and its
recorded Arb run, whereas the central-bin computation has a fresh MPFR
verification.  The new central pair has not yet received an independent
clean-room reproduction or peer review, and the argument is not Lean-verified.


The Price package has also been audited and rerun independently by the
\texttt{occisn} project \cite{Occisn2026}.  Other recent lower-bound work uses
agent-discovered convex relaxations \cite{KimPilanci2026} and Bochner/moment
constraints \cite{Zanghi2026}; these bounds are numerically weaker than the
Price certificate but provide useful alternative verification and relaxation
machinery.  On the upper side, public AI-guided construction searches include
TTT-Discover \cite{Yuksekgonul2026}.  We use none of these upper constructions
in the lower-bound proof.

\begin{theorem}\label{thm:main}
The minimum-overlap constant satisfies
\[
 \boxed{c_E>0.3805603}.
\]
\end{theorem}

\section{Continuous model and the positive-part dual}

Let $h:[0,2]\to[0,1]$ be measurable with $\int_0^2h=1$, put $g=1-h$
on $[0,2]$, and extend both functions by zero.  Define
\[
 p_h(t)=\int_{\R}h(x)g(x+t)\dd x,
 \qquad -2\le t\le2.
\]
Then $p_h\ge0$, $\int_{-2}^2p_h=1$, and the standard step-function reduction
gives
\[
 c_E\ge \inf_h\|p_h\|_\infty.
\]
We refer to \cite{White2022,Price2026} for a complete proof of the reduction
and for the ordinary moment and Fourier inequalities used by the certificate.

The following elementary dual principle is the mechanism used throughout.

\begin{lemma}[Positive-part dual]\label{lem:dual}
Suppose functions $\phi_j$ and constants $B_j$ satisfy
\[
 \int_{-2}^2p(t)\phi_j(t)\dd t\le B_j
\]
for every admissible profile in a fixed mean bin.  If $\lambda_j\ge0$ and
\[
 q(t)=1+\sum_j\lambda_j(B_j-\phi_j(t)),
\]
then
\[
 \|p\|_\infty\ge
 \left(\int_{-2}^2\max(q(t),0)\dd t\right)^{-1}.
\]
\end{lemma}

\begin{proof}
The assumed inequalities give $\int pq\ge1$.  Since $p\ge0$,
\[
 1\le\int pq\le\int p q_+
 \le\|p\|_\infty\int q_+.
\]
\end{proof}

\section{The Parseval-prefix inequality}

For $k\in\Z$ define
\[
 H_k=\int_0^2h(x)e^{ik\pi(x-1)}\dd x.
\]
At nonzero integer modes the transform of $g=1-h$ equals $-H_k$, because
\[
 \int_0^2e^{ik\pi(x-1)}\dd x=0.
\]
The Fourier transform of the overlap therefore gives
\begin{equation}\label{eq:integer-cos}
 C_k:=\int_{-2}^2p(t)\cos(k\pi t)\dd t=-|H_k|^2,
 \qquad k\ge1.
\end{equation}

\begin{lemma}[Parseval identity]\label{lem:parseval}
For every admissible $h$,
\[
 \boxed{p(0)+\sum_{k=1}^\infty |H_k|^2=\frac12}.
\]
Consequently, for every integer $K\ge1$,
\begin{equation}\label{eq:prefix}
 \boxed{-\sum_{k=1}^K\int_{-2}^2p(t)\cos(k\pi t)\dd t\le\frac12}.
\end{equation}
\end{lemma}

\begin{proof}
With the period-$2$ Fourier normalization, Parseval gives
\[
 1+2\sum_{k\ge1}|H_k|^2=2\int_0^2h(x)^2\dd x,
\]
because $H_0=\int_0^2h=1$.  Also
\[
 p(0)=\int_0^2h(x)(1-h(x))\dd x
 =1-\int_0^2h(x)^2\dd x.
\]
Combining the two identities proves the equality.  Since $p(0)\ge0$, the
finite prefix bound follows from \eqref{eq:integer-cos}.
\end{proof}

\begin{remark}
The force of \eqref{eq:prefix} is joint: the integer Fourier modes share one
energy budget.  Summing unrelated one-frequency inequalities loses this
coupling.
\end{remark}

\section{The central replacement certificate}

The pinned upstream certificate partitions the possible first moment into 172
bins.  Its binding symmetric pair consists of bins 85 and 86,
\[
 [-0.003125,0],\qquad [0,0.003125].
\]
The new certificate uses only even rows, so the same multipliers cover both
bins.  It contains the rows shown in Table~\ref{tab:rows}.

\begin{table}[ht]
\centering
\begin{tabular}{lr}
\toprule
Row type & Number \\
\midrule
Second-moment row & 1 \\
Ordinary cosine upper rows & 71 \\
Integer-$\pi$ cosine upper rows & 75 \\
Parseval-prefix rows, $K=191,195,200$ & 3 \\
\midrule
Total & 150 \\
\bottomrule
\end{tabular}
\caption{Rows in the central dual certificate.}
\label{tab:rows}
\end{table}

For the associated polynomial $q$, the directed-rounding checker proves
\begin{equation}\label{eq:central-D}
 \int_{-2}^2 q_+(t)\dd t
 \le 2.62770434388489007452190373282.
\end{equation}
At the theorem target,
\[
 \frac1{0.3805603}
 =2.6277044662830042965\ldots,
\]
so the certified $D$-side margin is larger than $1.22\times10^{-7}$.
Lemma~\ref{lem:dual} therefore proves the target on both central bins.

For the other 170 bins, the pinned upstream per-bin Arb report gives
\[
 D_I\le2.627538530873375790090272175878307131812586378826,
\]
which is below $1/0.3805603$ by more than $1.65\times10^{-4}$.  Thus every
outer certificate remains valid at the new target.

\begin{proof}[Proof of Theorem~\ref{thm:main}]
The new symmetric central certificate and the 170 pinned outer certificates cover
the full mean range $[-1,1]$.  Let
\[
 \begin{aligned}
 D_{\mathrm c}&:=2.62770434388489007452190373282,\\
 D_{\mathrm o}&:=2.627538530873375790090272175878307131812586378826,\\
 D_*&:=\max\{D_{\mathrm c},D_{\mathrm o}\}.
 \end{aligned}
\]
The certified bounds above give
$D_*<1/0.3805603$, uniformly over all 172 bins.  Lemma~\ref{lem:dual}
therefore yields
\[
 \|p_h\|_\infty\ge \frac1{D_*}>0.3805603
\]
for every admissible $h$, and hence
$\inf_h\|p_h\|_\infty\ge1/D_*>0.3805603$.  The standard step-function
reduction then gives $c_E>0.3805603$.
\end{proof}

\section{Directed-rounding verification}

The optimizer that produced the multipliers is not trusted.  The trusted input
is the decimal certificate and the analytic inequalities.  The C checker uses
96-bit MPFR endpoint arithmetic with directed rounding.  It performs Taylor
sign tests for $q$ on a dyadic subdivision of $[-2,2]$.  On certified-positive
components it uses the explicit antiderivative of the trigonometric polynomial;
on unresolved terminal cells it adds a conservative upper rectangle.

In this version, the SHA-256 digest of
\texttt{certificate/central\_certificate.json} is
\begin{center}
\small\texttt{8e085e786271759acb9063d82e84edcfdd3d73b6f556eecd39bc98c4f9911a8b}.
\end{center}

The recorded reference run reports
\[
 \begin{aligned}
 D_{\rm upper}&=2.62770434388489007452190373282,\\
 (1/C)_{\rm lower}&=2.62770446628300429655957281927,\\
 C_{\rm implied,lower}&=0.380560317726447488726080215466.
 \end{aligned}
\]
It visits $91{,}420$ cells and combines certified-positive cells into eight
components.  A separate embedding checker compares the embedded frequency
lists, multipliers, ordinary-cosine bounds, second-moment data, and target with
the corresponding entries in the JSON certificate.

The 170 outer bins are tied to the exact upstream commit and SHA-256 list.  The
repository includes a fetcher and a checker that recompute the maximum outer
$D_I$ from the pinned CSV report.  It does not locally rerun the upstream Arb
verification.  The trust boundary and reproduction commands are documented
with the source package.

\section{Discussion}

The improvement is numerically modest,
\[
 0.3805603-0.380554702762594012\ldots
 \approx5.60\times10^{-6},
\]
but it is a strict certified advance for a mature open problem.  More
importantly, it confirms that joint Fourier realizability information can
strengthen the compact profile-space dual.

A finite integer prefix uses only the radial energy at integer modes.  Natural
next steps are shifted Parseval lattices $((k+\theta)\pi)$, which couple cosine
and sine information, and exact multi-frequency bounded-density support cuts.
These may eventually be fitted to the contact structure of the best known
upper constructions.

\section*{Acknowledgments}
The author thanks Ethan Patrick White for the Fourier/convex framework and Liam
Price for publishing the 172-bin certificate and Arb verifier on which the
outer-bin part of this result is based.  The author thanks the maintainers of
the \texttt{occisn} project for their public independent audit of the upstream
certificate.

\section*{AI-assistance disclosure}
AI-assisted tools were used during exploratory research, drafting, and software
development.  They are not authors and their outputs are not treated as
mathematical evidence.  The named author is responsible for the statement,
proof, attribution, certificate, and released software.

\begin{thebibliography}{99}

\bibitem{Erdos1955}
P.~Erd\H{o}s,
\emph{Some remarks on number theory},
Riveon Lematematika \textbf{9} (1955), 45--48 (in Hebrew).

\bibitem{Moser1959}
L.~Moser,
\emph{On the minimal overlap problem of Erd\H{o}s},
Acta Arith. \textbf{5} (1959), no.~2, 117--119,
\href{https://doi.org/10.4064/aa-5-2-117-119}{doi:10.4064/aa-5-2-117-119}.

\bibitem{Haugland1996}
J.~K. Haugland,
\emph{Advances in the minimum overlap problem},
J. Number Theory \textbf{58} (1996), no.~1, 71--78,
\href{https://doi.org/10.1006/jnth.1996.0064}{doi:10.1006/jnth.1996.0064}.

\bibitem{Haugland2016}
J.~K. Haugland,
\emph{The minimum overlap problem revisited},
arXiv:1609.08000, 2016.

\bibitem{White2022}
E.~P. White,
\emph{Erd\H{o}s' minimum overlap problem},
arXiv:2201.05704, 2022.

\bibitem{Price2026}
L.~Price,
\emph{A certified lower bound for minimum overlap},
GitHub repository \texttt{Leeham06972452/erdos-36-lower-bound},
pinned commit \texttt{6bc610e40083ef61a40966dfb5d38612cabc4c5b}, 2026.


\bibitem{Occisn2026}
The \texttt{occisn} project,
\emph{A certified lower bound for Erd\H{o}s's minimum-overlap problem},
GitHub repository \texttt{occisn/erdos-36-certified-lower-bound}, 2026.

\bibitem{KimPilanci2026}
S.~Kim and M.~Pilanci,
\emph{AI-assisted discovery of convex relaxations via dual agents},
arXiv:2606.31182, 2026.

\bibitem{Zanghi2026}
B.~Zanghi,
\emph{Erd\H{o}s minimum overlap: Bochner and moment relaxations},
GitHub repository \texttt{bzanghi/erdos-minimum-overlap-bochner}, 2026.

\bibitem{Yuksekgonul2026}
M.~Yuksekgonul, D.~Koceja, X.~Li, F.~Bianchi, J.~McCaleb,
X.~Wang, J.~Kautz, Y.~Choi, J.~Zou, C.~Guestrin, and Y.~Sun,
\emph{Learning to discover at test time},
arXiv:2601.16175, 2026.

\bibitem{Johansson2017}
F.~Johansson,
\emph{Arb: efficient arbitrary-precision midpoint-radius interval arithmetic},
IEEE Trans. Comput. \textbf{66} (2017), no.~8, 1281--1292,
\href{https://doi.org/10.1109/TC.2017.2690633}{doi:10.1109/TC.2017.2690633},
arXiv:1611.02831.

\end{thebibliography}


\end{document}
